
\keyheading{第1题　在互相平行的直线下面打勾}
\begingroup\color{answerred}②画“$\checkmark$”。\par\endgroup
\keyheading{第2题　按要求画一画}
\begingroup\color{answerred}示例（答案不唯一）：\par\noindent\begin{center}{%
\begin{tikzpicture}[x=1bp,y=-1bp,color=black,draw=black,text=black]
\begin{scope}\clip (55.00,235.00) rectangle (553.00,561.00);
% 高清原题重排：PDF 21，原书 16。不含解答。
\hdhead{16}{相交与平行——信息窗2}
\hdsection{1}{141}{在互相平行的直线下面的（　　）里画“\(\checkmark\)”。}
\q{98}{186}{34}{13.5}{①}
\q{274}{186}{34}{13.5}{②}
\q{473}{186}{34}{13.5}{③}
\q{662}{186}{34}{13.5}{④}
\q{121}{247}{112}{13.5}{（　　）}
\q{310}{247}{112}{13.5}{（　　）}
\q{503}{247}{112}{13.5}{（　　）}
\q{675}{247}{112}{13.5}{（　　）}
\begin{scope}[scale=.6666667,draw=black,line width=.65pt]
\draw (128,205)--(178,205) (155,183)--(155,227);
\draw (305,210)--(377,183) (323,217)--(388,192.625);
\draw (516,210)--(555,185) (505,226)--(559,205);
\draw (695,187)--(735,187) (715,187)--(715,233);
\end{scope}
\hdsection{2}{305}{按要求画一画。}
\q{96}{348}{31}{13.5}{①}
\q{129}{349}{326}{13.5}{在点子图上任意画出两条平}
\q{129}{393}{326}{13.5}{行线。}
\q{472}{348}{31}{13.5}{②}
\q{505}{349}{340}{13.5}{在下面这条直线的左右两侧分}
\q{505}{393}{340}{13.5}{别画出一条与它平行的直线。}
\q{95}{581}{31}{13.5}{③}
\q{128}{582}{325}{13.5}{在方格纸上画一组垂线。}
\q{471}{581}{31}{13.5}{④}
\q{504}{582}{340}{13.5}{分别画出下面这条直线的一条}
\q{504}{626}{340}{13.5}{垂线和一条平行线。}
\begin{scope}[scale=.6666667,draw=black]
\foreach \i in {0,...,12}{\foreach \j in {0,...,7}{
  \fill ({186+16.75*\i},{440+111/7*\j}) circle (1.25);
  \fill ({563+203/12*\i},{440+111/7*\j}) circle (1.25);
}}
\draw[line width=.75pt] (655,536)--(707,440);
\foreach \i in {0,...,9}{
  \draw[hdgrid] ({177+188/9*\i},684)--({177+188/9*\i},813);
  \draw[hdgrid] ({569+192/9*\i},682)--({569+192/9*\i},812);
}
\foreach \j in {0,...,6}{
  \draw[hdgrid] (177,{684+129/6*\j})--(365,{684+129/6*\j});
  \draw[hdgrid] (569,{682+130/6*\j})--(761,{682+130/6*\j});
}
\draw[line width=.75pt] (569,747)--(761,747);
\end{scope}
\hdsection{3}{862}{用红笔描出两地之间最短的路线，并写出你的理由。}
\q{96}{915}{35}{13.5}{①}
\q{472}{916}{35}{13.5}{②}
\q{130}{1004}{80}{13.5}{甲港}
\q{360}{1004}{85}{13.5}{乙港}
\q{593}{914}{100}{13.5}{村庄}
\q{718}{1072}{100}{13.5}{公路}
\q{122}{1126}{87}{12.5}{理由：}
\q{503}{1126}{87}{12.5}{理由：}
\begin{scope}[scale=.6666667,draw=black,line width=.6pt]
\draw (185,1012) .. controls (188,982) and (198,991) .. (217,1002)
 .. controls (231,1007) and (240,1008) .. (245,998)
 .. controls (253,984) and (268,978) .. (274,985)
 .. controls (281,997) and (288,1017) .. (301,1000)
 .. controls (314,986) and (329,983) .. (335,992)
 .. controls (342,1000) and (350,1006) .. (356,1012);
\draw (185,1012) .. controls (235,1057) and (304,1056) .. (356,1012);
\draw (185,1012)--(356,1012);
\draw (523,1081)--(618,940)--(703,1081);
\draw (618,940)--(618,1081);
\draw (618,940) .. controls (610,959) and (613,975) .. (603,989)
 .. controls (596,999) and (589,1003) .. (596,1012)
 .. controls (586,1017) and (594,1024) .. (594,1030)
 .. controls (585,1043) and (592,1055) .. (582,1081);
\draw (517,1081) rectangle (710,1089);
\draw (618,1075)--(624,1075)--(624,1081);
\draw (183,1151)--(398,1151) (560,1151)--(806,1151);
\end{scope}
\q{400}{1127}{20}{12.5}{。}
\q{808}{1127}{20}{12.5}{。}

\begin{scope}[color=answerred,draw=answerred,text=answerred]\node[anchor=center,inner sep=0pt,outer sep=0pt,text=answerred,font=\normalsize] at (231.333,170.000) {\(\checkmark\)};
\draw[answerred,line width=.8pt] (124.000,303.333)--(258.000,303.333);
\draw[answerred,line width=.8pt] (124.000,346.000)--(258.000,346.000);
\draw[answerred,line width=.8pt] (402.000,357.333)--(436.667,293.333);
\draw[answerred,line width=.8pt] (474.000,357.333)--(508.667,293.333);
\draw[answerred,line width=.8pt] (118.000,499.333)--(243.333,499.333);
\draw[answerred,line width=.8pt] (173.333,456.000)--(173.333,542.000);
\draw[answerred,line width=.7pt] (178.667,499.333)--(178.667,494.000)--(173.333,494.000);
\draw[answerred,line width=.8pt] (379.333,469.333)--(506.667,469.333);
\draw[answerred,line width=.8pt] (450.000,454.667)--(450.000,541.333);
\draw[answerred,line width=.7pt] (455.333,498.000)--(455.333,492.667)--(450.000,492.667);
\draw[answerred,line width=1.3pt] (123.333,674.667)--(237.333,674.667);
\draw[answerred,line width=1.3pt] (412.000,626.667)--(412.000,720.667);
\ans{121.33}{751.333}{166}{10}{答：两点之间，线段最短。}
\ans{374.00}{751.333}{179}{10}{答：点到直线，垂线段最短。}
\end{scope}
\end{scope}\pgfresetboundingbox\path[use as bounding box] (55.00,235.00) rectangle (553.00,561.00);
\end{tikzpicture}}
\end{center}\par\endgroup
\keyheading{第3题　最短路线及理由}
\begingroup\color{answerred}\noindent（1）答：甲港与乙港之间的直线段；理由：两点之间，线段最短。\par
\noindent（2）答：村庄到公路的垂线段；理由：点到直线的所有连线中，垂线段最短。\par\par\smallskip\noindent\begin{center}\begin{adjustbox}{max width=.88\linewidth}\begin{tikzpicture}[draw=answerred,text=answerred]
\draw (0,0)--(3,0);\fill (0,0) circle (1.2pt);\fill (3,0) circle (1.2pt);
\node[left] at (0,0){甲港};\node[right] at (3,0){乙港};
\begin{scope}[xshift=6.5cm]
\draw (-1.4,0)--(1.4,0) node[right]{公路};\draw (0,0)--(0,1.8) node[above]{村庄};
\draw (0,.18)--(.18,.18)--(.18,0);
\end{scope}
\end{tikzpicture}\end{adjustbox}\end{center}\par\endgroup

\begingroup\color{answerred}\textbf{多解核对：}第2题的自由作图可有不同位置和方向；只要满足平行或垂直条件即正确。最短路线由指定点或两点确定，不因工具不同而改变。\par\endgroup
